\chapter{Analyse du P\&L et cycle de vie d'une stratégie}\label{ch:hf:pnl-analytics-and-the-strategy-lifecycle}

Une stratégie de market making a gagné moins ce mois-ci que le mois précédent. La fourchette s'est resserrée, le volume a augmenté, un concurrent est arrivé sur une plateforme, et
quelqu'un a modifié un paramètre le douze~; le rapport du matin doit dire lequel de ces facteurs en est responsable. Dans l'année simulée de ce chapitre, le mois 9 a rapporté
45\,600~\$ de moins que le mois 8~: la fourchette plus étroite a coûté 51\,500~\$, le volume plus élevé a rendu 48\,000~\$, le concurrent a pris 57\,400~\$ et le changement de paramètre
8\,100~\$, avec 23\,400~\$ de bruit dans l'autre sens. Estimée à partir des données propres de la stratégie, sans connaître ce qui avait été introduit, la part du concurrent
ressort à 53\,500~\$ et celle du paramètre à 8\,100~\$.

\section{Attribution quotidienne d'un book de market making}

Le P\&L quotidien d'un teneur de marché (livre~8, chapitre~1) se décompose en ce qu'il a gagné grâce à la fourchette, ce qu'il a perdu face aux flux informés, ce que son
inventaire a rapporté ou fait perdre, et ses frais et rabais (livre~7, chapitre~23, et chapitre~1 ici). Pour chaque plateforme et chaque instrument, la part espérée
est un produit~:
\[
\text{P\&L}=\underbrace{V}_{\text{volume du marché}}\times\underbrace{s}_{\text{notre part}}\times\bigl(\underbrace{\kappa\,h}_{\text{gain capté}}
-\underbrace{a}_{\text{sélection adverse}}\bigr),
\]
avec $h$ la demi-fourchette du marché, $\kappa$ la part de celle-ci que le teneur conserve (son taux de capture, net des frais et rabais) et $a$ la sélection adverse
par action, mesurée par les mark-outs. Les facteurs sont observables séparément~: $V$ et $h$ à partir du marché, $s$ à partir des exécutions du teneur, $\kappa$ et $a$ à partir
de ses prix d'exécution et de leurs mark-outs. Une variation du produit est une variation de l'un des facteurs, et le rôle de l'attribution est de dire lequel.

L'année du chapitre (\cref{fig:hf:pnl-analytics-and-the-strategy-lifecycle:daily})~: un teneur de marché sur deux plateformes, 10~\% de 20 millions d'actions par jour sur l'une et
12~\% de 15 millions sur l'autre, conservant 80~\% d'une demi-fourchette d'un cent et perdant 0.35 cent par action à la sélection adverse. Y ont été introduits~: la demi-fourchette
se resserre de 10~\% et le volume du marché augmente de 15~\% au mois 9~; un concurrent arrive sur la plateforme B le jour 175 et prend 30~\% de la part du teneur sur celle-ci et 15~\%
de son taux de capture~; un changement de paramètre, à partir du jour 182, augmente la part du teneur de 10~\% et sa sélection adverse de 0.10 cent par action les jours
où il est activé.

\begin{omfigure}
\begin{tikzpicture}
\begin{axis}[width=11cm,height=5.2cm,xlabel={\small jour de négociation},ylabel={\small P\&L quotidien (milliers de dollars)},xmin=0,xmax=252,ymin=0,ymax=30,
  tick label style={font=\footnotesize},grid=major,grid style={gray!25},table/col sep=comma,unbounded coords=jump,
  legend style={font=\scriptsize,at={(0.03,0.05)},anchor=south west},legend columns=2]
\fill[gray!12] (axis cs:168,0) rectangle (axis cs:189,30);
\addplot[gray,thin] table[x=day,y=pnl]{figdata/hft/28-pnl-analytics-and-the-strategy-lifecycle/daily.csv};
\addplot[omThm,thick] table[x=day,y=roll]{figdata/hft/28-pnl-analytics-and-the-strategy-lifecycle/daily.csv};
\draw[omDef,dashed] (axis cs:175,0) -- (axis cs:175,30);
\draw[omDef,dotted,thick] (axis cs:182,0) -- (axis cs:182,30);
\node[font=\scriptsize,anchor=north east,omDef] at (axis cs:174,29) {concurrent};
\node[font=\scriptsize,anchor=north west,omDef] at (axis cs:183,29) {paramètre};
\legend{quotidien,moyenne mobile sur 20 jours}
\end{axis}
\end{tikzpicture}

\omcaption{Une année simulée du P\&L quotidien d'un teneur de marché sur deux plateformes et de sa moyenne mobile sur 20 jours~; en grisé, le mois 9, lorsque la demi-fourchette se resserre et que le volume du marché
augmente~; en tirets, l'arrivée du concurrent sur la plateforme B (jour 175)~; en pointillés, le changement de paramètre, activé à partir du jour 182 sur des jours tirés au sort. Données~:
\texttt{hf\_attrib.daily}.}\label{fig:hf:pnl-analytics-and-the-strategy-lifecycle:daily}
\end{omfigure}

La variation d'un mois à l'autre est attribuée en modifiant les facteurs un à un, dans un ordre fixé, des valeurs du mois 8 à celles du mois 9
(\cref{lst:hf:pnl-analytics-and-the-strategy-lifecycle:change})~: la demi-fourchette, puis le volume du marché, puis le concurrent, puis le paramètre. La
fourchette et le volume sont mesurés à partir des données de marché~; l'effet du concurrent et celui du paramètre proviennent des deux sections suivantes. Ce qui reste relève du bruit
et de l'approximation liée à l'usage de moyennes mensuelles.

\section{Changements de paramètres comme expériences}

\begin{definition}[Contrôle des changements de paramètres]\label{def:hf:pnl-analytics-and-the-strategy-lifecycle:control}
Le \emph{contrôle des changements de paramètres}\index{contrôle des changements de paramètres} est la discipline qui consiste à traiter chaque modification des paramètres d'une stratégie en production comme une
expérience enregistrée, revue et réversible~: qui a changé quoi, quand et pourquoi, la modification étant appliquée à une partie aléatoire de l'activité de la stratégie
jusqu'à ce que son effet soit mesuré.
\end{definition}

Un paramètre modifié le douze pour tout ce qui suit est confondu avec tout ce qui s'est produit après le douze. Appliqué
plutôt sur des jours tirés au sort (livre~7, chapitre~21~; le \texttt{firm.abtest} du livre~7 les attribue par un hachage), son effet est estimé en comparant
les jours traités et les jours non traités des mêmes semaines. Sur les 70 jours à partir du jour 182, 33 ont été traités~: le paramètre a réduit le P\&L quotidien de 2\,670~\$ (erreur type
1\,070~\$)~; un ajustement sur le volume quotidien du marché (CUPED) donne 2\,550~\$ (1\,070~\$). Les facteurs sont plus nets que le P\&L~: la part des jours traités
a été exactement supérieure de 10~\% et leur sélection adverse de 0.10 cent par action, parce que les deux sont mesurées sur chaque exécution. Le paramètre achète du volume à un
prix en toxicité qui dépasse sa valeur~; il faut le faire revenir en arrière.

\section{Érosion de la capture et concurrence}

\begin{definition}[Érosion de la capture]\label{def:hf:pnl-analytics-and-the-strategy-lifecycle:decay}
L'\emph{érosion de la capture}\index{erosion de la capture@érosion de la capture} est la baisse, au fil du temps, de la part de la fourchette que conserve une stratégie de market making, ou de la part
des flux qu'elle traite, à mesure que des concurrents la copient, que les plateformes évoluent ou que les flux qu'elle servait se détournent.
\end{definition}

Un concurrent se manifeste d'abord dans la part du teneur et dans son taux de capture sur la plateforme qu'il rejoint. Un test CUSUM unilatéral (livre~7, chapitre~13) appliqué
à la part quotidienne du teneur sur la plateforme B, par rapport à sa cible de 12~\% avec une marge d'un demi-point de pourcentage, franchit son seuil le jour 176, un jour après l'arrivée
(\cref{fig:hf:pnl-analytics-and-the-strategy-lifecycle:cusum})~; le même test appliqué au taux de capture franchit le seuil le même jour. En comparant les jours précédant
le point de rupture aux jours non traités qui le suivent, on estime l'effet du concurrent~: 29.9~\% de la part et 14.9~\% du taux de capture, contre les
30~\% et 15~\% introduits.

\begin{omfigure}
\begin{tikzpicture}
\begin{axis}[width=11cm,height=4.8cm,axis y line*=right,axis x line=none,xmin=150,xmax=199,ymin=0,ymax=1.2,
  ylabel={\small statistique CUSUM},tick label style={font=\footnotesize},table/col sep=comma]
\addplot[omDef,thick,dashed] table[x=day,y=cusum]{figdata/hft/28-pnl-analytics-and-the-strategy-lifecycle/cusum.csv};
\label{plot:hf:cusumstat}
\draw[gray,dotted] (axis cs:150,0.05) -- (axis cs:199,0.05);
\end{axis}
\begin{axis}[width=11cm,height=4.8cm,axis y line*=left,xmin=150,xmax=199,xlabel={\small jour de négociation},ylabel={\small part sur la plateforme B (\%)},
  ymin=0,ymax=15,tick label style={font=\footnotesize},grid=major,grid style={gray!25},table/col sep=comma,
  legend style={font=\scriptsize,at={(0.03,0.05)},anchor=south west}]
\addplot[omThm,thick,const plot] table[x=day,y=share]{figdata/hft/28-pnl-analytics-and-the-strategy-lifecycle/cusum.csv};
\addlegendentry{part sur la plateforme B (gauche)}
\addlegendimage{/pgfplots/refstyle=plot:hf:cusumstat}\addlegendentry{statistique CUSUM (droite)}
\end{axis}
\end{tikzpicture}

\omcaption{Détection du concurrent~: la part quotidienne du teneur dans le volume de la plateforme B et la statistique CUSUM unilatérale de son écart à 12~\% (marge de
0.5 point, seuil 0.05, pointillés), jours 150 à 199~; la part chute le jour 175 et la statistique franchit le seuil le jour 176~; les jours traités après le jour 182
montrent les 10~\% de part en plus dus au paramètre. Données~: \texttt{hf\_attrib.cusum\_path}.}\label{fig:hf:pnl-analytics-and-the-strategy-lifecycle:cusum}
\end{omfigure}

Avec ces deux estimations, le mois 9 s'attribue comme dans la \cref{fig:hf:pnl-analytics-and-the-strategy-lifecycle:month9}~: sur les 45\,600~\$ de baisse, l'attribution estimée
donne à la fourchette $-\$51\,500$, au volume $+\$48\,000$, au concurrent $-\$53\,500$ (introduit~: $-\$57\,400$) et au paramètre $-\$8\,100$ (introduit~:
$-\$8\,100$)~; le résidu, 19\,600~\$, est le bruit du mois sur l'inventaire, le volume et les fourchettes. L'effet estimé du concurrent est plus faible que l'effet
introduit parce que le point de rupture est détecté avec un jour de retard.

\begin{omfigure}
\begin{tikzpicture}
\begin{axis}[width=11cm,height=5cm,ybar,bar width=10pt,xmin=-0.6,xmax=4.6,xtick={0,1,2,3,4},
  xticklabels={fourchette,volume,concurrent,paramètre,résidu},ylabel={\small mois 9 moins mois 8 (milliers de dollars)},ymin=-70,ymax=60,
  tick label style={font=\footnotesize},grid=major,grid style={gray!25},table/col sep=comma,
  legend style={font=\scriptsize,at={(0.97,0.03)},anchor=south east}]
\addplot[fill=gray!40,draw=gray] table[x=i,y=truth]{figdata/hft/28-pnl-analytics-and-the-strategy-lifecycle/month9.csv};
\addplot[fill=omThm!60,draw=omThm] table[x=i,y=estimate]{figdata/hft/28-pnl-analytics-and-the-strategy-lifecycle/month9.csv};
\draw[black] (axis cs:-0.6,0) -- (axis cs:4.6,0);
\legend{introduit,estimé}
\end{axis}
\end{tikzpicture}

\omcaption{La variation mois sur mois du P\&L attribuée à la demi-fourchette, au volume de marché, au concurrent de la plateforme B et au changement de paramètre, avec
les valeurs introduites et les valeurs estimées à partir des données de la stratégie (le point de rupture et l'expérience)~; le résidu est ce que
n'explique aucun de ces facteurs. Données~: \texttt{hf\_attrib.month9}.}\label{fig:hf:pnl-analytics-and-the-strategy-lifecycle:month9}
\end{omfigure}

L'ordre d'attribution compte~: l'effet du concurrent, évalué avec la fourchette plus étroite du mois~9, est plus faible que celui évalué avec celle du mois~8.
Un ordre fixe et documenté, ou la moyenne sur les ordres (une attribution de Shapley), fait partie de la définition du rapport~; il en va de même de la présentation du résidu plutôt que de sa dissimulation.
La décomposition de Menkveld pour un teneur de marché à haute fréquence (chapitre~1~: 1.55~\euro{} par transaction sur la fourchette nette de frais, 0.68~\euro{} perdus sur les positions) repose
sur la même idée, au niveau de l'ensemble du livre d'un établissement.

\section{Quand retirer une stratégie}

\begin{definition}[Cycle de vie d'une stratégie]\label{def:hf:pnl-analytics-and-the-strategy-lifecycle:lifecycle}
Le \emph{cycle de vie d'une stratégie}\index{cycle de vie d'une stratégie} est la suite des étapes qu'une stratégie de trading traverse, depuis la recherche, la revue et un lancement progressif,
en passant par la production surveillée et les modifications menées sous forme d'expériences, jusqu'à la réduction et au retrait lorsque son edge mesuré ne couvre plus ses coûts, chaque
étape ayant ses critères fixés à l'avance.
\end{definition}

Une règle de retrait est un critère d'abandon (livre~7, chapitre~1) appliqué à une stratégie entière~: on arrête lorsque l'edge, mesuré sur une fenêtre
assez longue pour être fiable, passe sous ce que coûte la stratégie (personnel, technologie, capital). Dans l'année simulée, le teneur de marché gagnait 16\,600~\$ par jour avant le mois~8
; après l'arrivée du concurrent et avec la demi-fourchette réduite de moitié, 12\,200~\$. Une règle qui retire la stratégie lorsque sa moyenne sur 60~jours passe sous un
coût de 12\,000~\$ par jour se déclenche le jour~239~; à 13\,000~\$ par jour, le jour~221~; à 10\,000~\$, jamais. La règle est une décision à prendre avant la dégradation, et non
après~: le coût est connu à l'avance, la fenêtre fixe la vitesse à laquelle la règle peut agir et la fréquence de ses erreurs.

\section{Tutoriel~: lequel de ces facteurs l'a causé}\label{tut:hf:pnl-analytics-and-the-strategy-lifecycle:tut}

\begin{tutorial}
\textbf{Objectif.} Attribuer la variation mensuelle du P\&L d'un teneur de marché à ses causes, les causes étant estimées à partir des données propres de la stratégie.
\textbf{État final~:} les quatre chiffres et les nombres du texte.
\begin{tutsteps}
\item \textbf{L'année} avec des événements introduits (\texttt{firm.mmattrib.Year}) et son attribution quotidienne.
\item \textbf{Le concurrent}~: CUSUM sur la part de la plateforme B et sur le ratio de capture, puis l'effet avant et après le point de rupture.
\item \textbf{Le paramètre}~: l'expérience randomisée (\texttt{firm.mmattrib.experiment} sur \texttt{firm.abtest} du livre~7).
\omcode{code/firm/mmattrib/firm_mmattrib.py}{83}{96}{Jours traités contre jours non traités~: l'effet sur le P\&L quotidien, avec CUPED, et sur la part et la sélection adverse.}\label{lst:hf:pnl-analytics-and-the-strategy-lifecycle:exp}
\item \textbf{L'attribution}, facteur par facteur.
\omcode{code/firm/mmattrib/firm_mmattrib.py}{119}{135}{Modifier la demi-fourchette, puis le volume, puis le concurrent, puis le paramètre~; ce qui reste est le résidu.}\label{lst:hf:pnl-analytics-and-the-strategy-lifecycle:change}
\end{tutsteps}
\textbf{Pour aller plus loin.} Attribuer par moyenne de Shapley sur les ordres~; ajouter une troisième plateforme où le concurrent arrive aussi plus tard~; appliquer
l'expérience à la moitié des instruments au lieu de la moitié des jours.
\end{tutorial}

\section{Construction~: le moteur d'attribution}\label{bld:hf:pnl-analytics-and-the-strategy-lifecycle:mmattrib}

\begin{build}
\bfield{Rôle} Attribuer le P\&L d'un teneur de marché à ses facteurs, chaque jour et mois après mois, détecter les changements de capture, lancer les changements de paramètres
sous forme d'expériences et appliquer une règle de retrait.
\bfield{Interface} \texttt{Year(seed)} et ses paramètres d'événements introduits, \texttt{attribution(year)}, \texttt{cusum\_down(x, target, k, h)},
\texttt{experiment(year, first, last)}, \texttt{month\_change(year, m0, m1, est)}, \texttt{retire\_day(pnl, window, floor)}. Construit sur \texttt{firm.abtest} du livre~7.
\bfield{Règles} Dollars et centimes par action~; un ordre d'attribution fixe~; le résidu présenté.
\bfield{Tests d'acceptation} \texttt{code/firm/mmattrib/tests/}~: les parties s'additionnent au total et le ratio de capture observé égale celui introduit~; le
concurrent réduit la part de la plateforme B et l'expérience débute le jour prévu~; CUSUM se déclenche le lendemain d'un saut et jamais sans saut~; l'expérience
retrouve la part et la sélection adverse introduites~; sans bruit, l'attribution laisse moins de 5~\% inexpliqué~; la règle de retrait se déclenche après une baisse et
non avant.
\bfield{Pour aller plus loin} Attribution de Shapley~; attribution par instrument et par plateforme à la fois~; tests séquentiels pour l'expérience.
\end{build}

\begin{omsources}
\item A.~J.~Menkveld, High frequency trading and the new market makers, \emph{Journal of Financial Markets} 16(4), 2013, 712--740.
\item livre~7, chapitres~1, 13 et 21 (critères d'abandon, tests CUSUM, expériences).
\end{omsources}

\section{Exercices}

\begin{exercise}[$\star$]\label{exo:hf:pnl-analytics-and-the-strategy-lifecycle:1}
Une plateforme négocie 15 millions d'actions par jour~; la part du teneur de marché est de 12~\%, il conserve 80~\% d'une demi-fourchette d'un centime et perd 0.35 centime par action. Que
gagne-t-il par jour sur cette plateforme~?
\end{exercise}

\begin{exercise}[$\star$]\label{exo:hf:pnl-analytics-and-the-strategy-lifecycle:2}
Le concurrent ramène la part à 8.4~\% et le ratio de capture à 68~\%. Que gagne maintenant la plateforme, à fourchette égale~?
\end{exercise}

\begin{exercise}[$\star$]\label{exo:hf:pnl-analytics-and-the-strategy-lifecycle:3}
Le paramètre augmente la part de 10~\% et la sélection adverse de 0.10 centime. Est-il rentable sur la plateforme de l'exemple, avant l'arrivée du concurrent~?
\end{exercise}

\begin{exercise}[$\star\star$]\label{exo:hf:pnl-analytics-and-the-strategy-lifecycle:4}
Pourquoi un paramètre modifié pour tous les jours à partir du douze est-il difficile à évaluer~?
\end{exercise}

\begin{exercise}[$\star\star$]\label{exo:hf:pnl-analytics-and-the-strategy-lifecycle:5}
Pourquoi l'ordre d'attribution modifie-t-il la part de la baisse imputée au concurrent~?
\end{exercise}

\begin{exercise}[$\star\star$]\label{exo:hf:pnl-analytics-and-the-strategy-lifecycle:6}
Pourquoi la part et la sélection adverse sont-elles de meilleures mesures de l'effet du paramètre que le P\&L quotidien~?
\end{exercise}

\begin{exercise}[$\star\star\star$]\label{exo:hf:pnl-analytics-and-the-strategy-lifecycle:7}
\emph{Programmation.} Lancer l'attribution sur les graines 2 et 3. Quelle est l'ampleur du résidu, et à quel point les effets estimés du concurrent et du paramètre
sont-ils proches de ceux introduits~?
\end{exercise}

\begin{exercise}[$\star\star\star$]\label{exo:hf:pnl-analytics-and-the-strategy-lifecycle:8}
\emph{Trouver la faille.} «~Le P\&L a baissé après que nous avons modifié le paramètre le douze, donc c'est le changement de paramètre qui en est la cause~; il faut le rétablir.~»
\end{exercise}

\section{Problème~: lequel de ces facteurs l'a causé}

\begin{problem}[{Problème du week-end --- lequel de ces facteurs l'a causé}]\label{pb:hf:pnl-analytics-and-the-strategy-lifecycle:1}
Le rapport du matin d'un teneur de marché doit expliquer pourquoi le mois~9 a rapporté moins que le mois~8.

\medskip\noindent\textbf{Partie I --- Attribution.}
\begin{enumerate}
\item Écrire le P\&L quotidien espéré comme un produit de facteurs et indiquer où chacun est mesuré.
\item Décrire l'année simulée et ses événements introduits.
\item Décrire l'attribution séquentielle.
\item Que montre la décomposition de Menkveld au niveau d'un établissement~?
\end{enumerate}

\medskip\noindent\textbf{Partie II --- Estimer les causes.}
\begin{enumerate}[resume]
\item Définir le \omterm{def:hf:pnl-analytics-and-the-strategy-lifecycle:control}{contrôle des changements de paramètres}.
\item Donner les estimations de l'expérience, avec et sans CUPED.
\item Définir l'\omterm{def:hf:pnl-analytics-and-the-strategy-lifecycle:decay}{érosion de la capture} (capture decay) et décrire la détection par le test CUSUM.
\item Donner l'effet concurrent estimé et le comparer à l'effet introduit.
\end{enumerate}

\medskip\noindent\textbf{Partie III --- L'attribution.}
\begin{enumerate}[resume]
\item Donner l'attribution introduite et l'attribution estimée du mois~9.
\item Pourquoi l'effet concurrent estimé est-il plus petit~?
\item Quel est le résidu~?
\item Pourquoi l'ordre de l'attribution doit-il être fixé et documenté~?
\end{enumerate}

\medskip\noindent\textbf{Partie IV --- Le verdict.}
\begin{enumerate}[resume]
\item Énoncer le \emph{résultat nommé}~: l'attribution de la variation mensuelle du P\&L à la fourchette, au volume, à la concurrence et au changement de paramètre,
face à la vérité introduite.
\item Faut-il conserver le paramètre~?
\item Définir le \omterm{def:hf:pnl-analytics-and-the-strategy-lifecycle:lifecycle}{cycle de vie d'une stratégie}.
\item Quand la règle de retrait se déclenche-t-elle, et de quoi dépend-elle~?
\item Que doit montrer le rapport chaque matin~?
\item En quoi une attribution de Shapley différerait-elle~?
\item Que ferait le concurrent à une stratégie incapable de randomiser ses changements~?
\item En une phrase~: que rend possible l'attribution~?
\end{enumerate}
\end{problem}

\section{Questions d'entretien}

\begin{interviewq}[$\star$ \iqroles{trader}]\label{iq:hf:pnl-analytics-and-the-strategy-lifecycle:1}
Le P\&L de votre stratégie a été divisé par deux cette semaine. Que regardez-vous en premier~?
\end{interviewq}

\begin{interviewq}[$\star\star$ \iqroles{chercheur}]\label{iq:hf:pnl-analytics-and-the-strategy-lifecycle:2}
Comment détecteriez-vous, à partir de vos propres données, qu'un concurrent a rejoint votre plateforme~?
\end{interviewq}

\begin{interviewq}[$\star\star$ \iqroles{développeur}]\label{iq:hf:pnl-analytics-and-the-strategy-lifecycle:3}
Concevez le système qui enregistre chaque changement de paramètre et son affectation à l'expérience, et qui produit le rapport d'attribution pour 7~h.
\end{interviewq}

\begin{interviewq}[$\star\star$ \iqroles{risque}]\label{iq:hf:pnl-analytics-and-the-strategy-lifecycle:4}
Un trader veut modifier un paramètre en production maintenant, et non dans le cadre d'une expérience. Que exigez-vous~?
\end{interviewq}

\begin{interviewq}[$\star\star$ \iqroles{chercheur}]\label{iq:hf:pnl-analytics-and-the-strategy-lifecycle:5}
Combien de jours une expérience doit-elle durer pour détecter un effet de 2\,000~\$ par jour, avec un écart type quotidien de 4\,500~\$~?
\end{interviewq}

\begin{interviewq}[$\star\star\star$ \iqroles{chercheur}]\label{iq:hf:pnl-analytics-and-the-strategy-lifecycle:6}
Montrez qu'une attribution séquentielle d'un produit de deux facteurs affecte le terme d'interaction au facteur modifié en second.
\end{interviewq}
